% ============================================================
\part{Entwicklungsumgebung aufbauen}
% ============================================================

\chapter{Wir wollen ein Projekt übersichtlich anlegen}

\kapitelziel{
Du kennst die vollständige Verzeichnisstruktur von
robotik-uno-q und verstehst, warum jeder Ordner
seinen Platz hat.
}
\kapitelstruktur

\section{Situation}

Das Projekt wächst. \texttt{stufe2\_ros2/} hat jetzt
\texttt{src/}, \texttt{config/}, \texttt{logs/} und bald
kommen Tests, Dokumentation und Launch-Dateien dazu.
Gleichzeitig liegt \texttt{stufe1\_mcu/} mit Arduino-Code
im selben Repo. Und \texttt{docs/} mit Schaltplänen.

Wer das Projekt öffnet, soll in 30 Sekunden verstehen,
was wo liegt -- ohne zu fragen.

\begin{aufgabeBox}
Vervollständige die \texttt{stufe2\_ros2/}-Struktur auf
ROS2-kompatiblen Standard und dokumentiere die Projektstruktur
in \texttt{docs/}.
\end{aufgabeBox}

\section{Die vollständige Projektstruktur}

\begin{lstlisting}[style=text]
robotik-uno-q/
├── README.md                  Was ist das Projekt?
├── .gitignore                 Was ignoriert Git?
│
├── stufe1_mcu/                Arduino C++ (laeuft schon)
│   ├── differentialantrieb/
│   ├── hinderniserkennung/
│   └── linienfolger/
│
├── stufe2_ros2/               Python + ROS2 (wir bauen das)
│   ├── src/
│   │   └── sensor_bridge.py
│   ├── config/
│   │   └── sensor_config.yaml
│   ├── launch/                (in Band 2: ros2 launch)
│   ├── tests/                 (spaeter: pytest)
│   └── logs/                  (in .gitignore!)
│
├── stufe3_wifi/               WiFi-Kommunikation (Band 3)
│
├── docs/
│   ├── Verdrahtung.md         Schaltplan Stufe 1
│   ├── architektur.md         (heute anlegen)
│   └── pics/
│
├── notebooks/                 Jupyter (Kinematik, Kalibrierung)
│
├── atopile/                   KiCad PCB
└── scripts/
    └── start_roboter.sh
\end{lstlisting}

\begin{loesungBox}
\begin{lstlisting}[style=bash]
cd ~/robotik-uno-q

# Fehlende Ordner erganzen
mkdir -p stufe2_ros2/launch stufe2_ros2/tests

# Architektur-Dokument anlegen
cat > docs/architektur.md << 'MDEOF'
# Systemarchitektur robotik-uno-q

## Stufen

| Stufe | Technologie | Status |
|-------|------------|--------|
| 1     | Arduino C++ (L298N, HC-SR04) | fertig |
| 2     | Python + ROS2 | in Entwicklung |
| 3     | WiFi / MQTT | geplant |

## Kommunikation

stufe1_mcu  <-- Serial /dev/ttyUSB0 -->  stufe2_ros2
stufe2_ros2  <-- ROS2 Topics -->  Navigation + Motor
MDEOF

echo "Struktur komplett."
tree robotik-uno-q/   # apt install tree
\end{lstlisting}
\end{loesungBox}

\begin{projektBox}
\textbf{Bisher:} Teile der Struktur gewachsen.\\
\textbf{Heute:} Vollständige, dokumentierte Projektstruktur.
Jeder Ordner hat einen Zweck.\\
\textbf{Nächstes Kapitel:} VS Code als komfortable IDE einrichten.
\end{projektBox}

\begin{merksatzBox}
Eine gute Projektstruktur ist die Einladungskarte für Kollegen.
Wer das Repo klont, soll in 30 Sekunden wissen: wo ist der Code?
Wo sind Tests? Wo ist die Dokumentation?
\texttt{README.md} beantwortet \emph{Was ist das?}
\texttt{docs/architektur.md} beantwortet \emph{Wie funktioniert es?}
\end{merksatzBox}

\begin{robotikBox}
Ein ROS2-Paket erzwingt eine Struktur: \texttt{package.xml},
\texttt{setup.py}, \texttt{resource/}, \texttt{launch/}.
Die Ordner \texttt{stufe2\_ros2/src/}, \texttt{config/} und
\texttt{launch/} entsprechen exakt dieser Konvention.
Wenn Band~2 das Paket formalisiert, ist fast nichts zu verschieben.
\end{robotikBox}

% ── Kapitel 12 ────────────────────────────────────────────────
\chapter{Wir wollen komfortabel programmieren}

\kapitelziel{
Du richtest VS Code für robotik-uno-q ein: mit
Arduino-Unterstützung, Python-Linting und direktem
Remote-Zugriff auf den Raspberry Pi.
}
\kapitelstruktur

\section{Situation}

Du entwickelst gleichzeitig Arduino-Code in \texttt{stufe1\_mcu/}
und Python in \texttt{stufe2\_ros2/}. Mit \texttt{nano} allein
ist das mühsam: keine Autovervollständigung, kein Syntax-Highlighting
für \texttt{.ino}-Dateien, kein direkter Upload.

\begin{aufgabeBox}
Richte VS Code so ein, dass du in einer IDE sowohl
Arduino- als auch Python-Code für robotik-uno-q entwickelst --
und bei Bedarf direkt auf dem Raspberry Pi.
\end{aufgabeBox}

\section{VS Code einrichten}

\begin{lstlisting}[style=bash]
# VS Code installieren (Ubuntu/Debian)
sudo apt install code   # oder von code.visualstudio.com

# Extensions installieren (Strg+Shift+X):
# - Python          (ms-python.python)
# - Arduino         (vsciot-vscode.vscode-arduino)
# - YAML            (redhat.vscode-yaml)
# - Remote - SSH    (ms-vscode-remote.remote-ssh)
# - GitLens         (eamodio.gitlens)
\end{lstlisting}

\begin{lstlisting}[style=json]
// .vscode/settings.json im Projektordner
{
  "python.defaultInterpreterPath": "/usr/bin/python3",
  "arduino.path": "/usr/share/arduino",
  "arduino.additionalUrls": [
    "https://raw.githubusercontent.com/stm32duino/BoardManagerFiles/main/package_stmicroelectronics_index.json"
  ],
  "yaml.schemas": {
    "file://config/sensor_config.yaml": "stufe2_ros2/config/*.yaml"
  }
}
\end{lstlisting}

\section{Remote-SSH auf den Raspberry Pi}

\begin{lstlisting}[style=bash]
# SSH-Key generieren (einmalig)
ssh-keygen -t ed25519 -C "robotik-uno-q"
ssh-copy-id ros@192.168.1.100   # IP des Raspberry Pi

# ~/.ssh/config
Host roboter-pi
    HostName 192.168.1.100
    User ros
    IdentityFile ~/.ssh/id_ed25519
\end{lstlisting}

\begin{lstlisting}[style=text]
Dann in VS Code:
Strg+Shift+P -> "Remote-SSH: Connect to Host"
-> "roboter-pi"
-> Ordner oeffnen: ~/robotik-uno-q
\end{lstlisting}

\begin{loesungBox}
Nach der Einrichtung:
\begin{itemize}
  \item Lokaler Rechner: VS Code öffnet das Projekt.
  \item Arduino-Code: \texttt{.ino}-Dateien mit Syntax-Highlighting.
  \item Python-Code: Autovervollständigung und Linting.
  \item Raspberry Pi: Remote-SSH -- Code läuft auf dem Pi, Editor ist lokal.
  \item Git: Änderungen sind sofort im Source-Control-Tab sichtbar.
\end{itemize}
\end{loesungBox}

\begin{projektBox}
\textbf{Bisher:} Vollständige Projektstruktur.\\
\textbf{Heute:} VS Code mit Arduino- und Python-Support,
Remote-SSH für den Raspberry Pi.\\
\textbf{Nächstes Kapitel:} Git-History des Projekts weiterführen.
\end{projektBox}

\begin{merksatzBox}
Remote-SSH ist der Industriestandard für Embedded-Entwicklung.
Der Code läuft auf dem Roboter, du schreibst ihn am Laptop.
Kein manuelles Kopieren, kein Dateiübertragen -- VS Code
synchronisiert alles über SSH.
\end{merksatzBox}

\begin{robotikBox}
VS Code hat eine offizielle ROS2-Extension, die den Node-Graph
visualisiert, Topics anzeigt und direkt \texttt{ros2 run}
aus der IDE startet. Wenn Band~2 beginnt, ist VS Code
bereits als IDE eingerichtet -- nur die ROS2-Extension kommt noch dazu.
\end{robotikBox}
